%%%%%%%%%%%%%%%%%%%%%%%%%%%%%%%%%%%%%%%%%%%%%%%%%%%%%%%%%%
\begin{frame}
  \begin{center}
    {\Large Introduction to Pandas}
  \end{center}
\end{frame}

%%%%%%%%%%%%%%%%%%%%%%%%%%%%%%%%%%%%%%%%%%%%%%%%%%%%%%%%%%
\begin{frame}[fragile]\frametitle{Pandas}
\textbf{\textit{pandas}} introduces two new data structures to Python - \textbf{Series} and \textbf{DataFrame}, both of which are built on top of NumPy. Throughout this session, lines starting with \texttt{\textgreater\textgreater\textgreater} (or \texttt{...} for a continued line) are what you type at the Python prompt; everything else is the output pandas prints back.
\begin{lstlisting}
>>> import pandas as pd
>>> import numpy as np
>>> import matplotlib.pyplot as plt
>>> pd.set_option('display.max_columns', 50)
\end{lstlisting}
\end{frame}

%%%%%%%%%%%%%%%%%%%%%%%%%%%%%%%%%%%%%%%%%%%%%%%%%%%%%%%%%%
\begin{frame}[fragile]\frametitle{Series}
Series is a one-dimensional labeled array capable of holding any data type (integers, strings, floating point numbers, Python objects, etc.). The axis labels are collectively referred to as the index. The basic method to create a Series is to call:

 Here, data can be many different things:
\begin{itemize}
\item a Python \texttt{dict}
\item an \texttt{ndarray}
\item a scalar value (like 5)
\end{itemize}

\begin{lstlisting}
>>> s = pd.Series(data, index=index)
\end{lstlisting}
\end{frame}

%%%%%%%%%%%%%%%%%%%%%%%%%%%%%%%%%%%%%%%%%%%%%%%%%%%%%%%%%%%%%%%%%%%%%%%%
\begin{frame}[fragile]\frametitle{Creating a Series from a List}
\begin{itemize}
\item A Series is a one-dimensional object similar to an array, list, or column in a table.
\item It will assign a labeled index to each item in the Series. \item By default, each item will receive an index label from 0 to N, where N is the length of the Series minus one.
\end{itemize}
\begin{lstlisting}
>>> # create a Series with an arbitrary list
>>> s = pd.Series([7, 'Heisenberg', 3.14, -1789710578, 'Happy Eating!'])
>>> s
\end{lstlisting}
\end{frame}

%%%%%%%%%%%%%%%%%%%%%%%%%%%%%%%%%%%%%%%%%%%%%%%%%%%%%%%%%%%%%%%%%%%%%%%%
\begin{frame}[fragile]\frametitle{Series}
\textbf{Output from Previous Slide}
\begin{lstlisting}
0                7
1       Heisenberg
2             3.14
3      -1789710578
4    Happy Eating!
dtype: object
\end{lstlisting}
\end{frame}

%%%%%%%%%%%%%%%%%%%%%%%%%%%%%%%%%%%%%%%%%%%%%%%%%%%%%%%%%%%%%%%%%%%%%%%%
\begin{frame}[fragile]\frametitle{Series with a Custom Index}
Alternatively, you can specify an index to use when creating the Series.
\begin{lstlisting}
>>> s = pd.Series([7, 'Heisenberg', 3.14, -1789710578,
...    'Happy Eating!'],
... index=['A', 'Z', 'C', 'Y', 'E'])
>>> s

A                7
Z       Heisenberg
C             3.14
Y      -1789710578
E    Happy Eating!
dtype: object
\end{lstlisting}
\end{frame}

%%%%%%%%%%%%%%%%%%%%%%%%%%%%%%%%%%%%%%%%%%%%%%%%%%%%%%%%%%%%%%%%%%%%%%%%
\begin{frame}[fragile]\frametitle{Series}
The Series constructor can convert a dictonary as well, using the keys of the dictionary as its index.
\begin{lstlisting}
>>> d = {'Chicago': 1000, 'New York': 1300, 'Portland': 900, 'San Francisco': 1100,'Austin': 450, 'Boston': None}
>>> cities = pd.Series(d)
>>> cities
Austin            450
Boston            NaN
Chicago          1000
New York         1300
Portland          900
San Francisco    1100
dtype: float64
\end{lstlisting}
\end{frame}

%%%%%%%%%%%%%%%%%%%%%%%%%%%%%%%%%%%%%%%%%%%%%%%%%%%%%%%%%%%%%%%%%%%%%%%%
\begin{frame}[fragile]\frametitle{Series}
You can use the index to select specific items from the Series
\begin{lstlisting}
>>> cities['Chicago']
1000.0
\end{lstlisting}
\end{frame}

%%%%%%%%%%%%%%%%%%%%%%%%%%%%%%%%%%%%%%%%%%%%%%%%%%%%%%%%%%%%%%%%%%%%%%%%
\begin{frame}[fragile]\frametitle{Series}
\begin{lstlisting}
>>> cities[['Chicago', 'Portland', 'San Francisco']]
Chicago          1000
Portland          900
San Francisco    1100
dtype: float64
\end{lstlisting}
\end{frame}

%%%%%%%%%%%%%%%%%%%%%%%%%%%%%%%%%%%%%%%%%%%%%%%%%%%%%%%%%%%%%%%%%%%%%%%%
\begin{frame}[fragile]\frametitle{Series}
You can use \textbf{\textit{boolean indexing}} for selection.
\begin{lstlisting}
>>> cities[cities < 1000]
Austin      450
Portland    900
dtype: float64
\end{lstlisting}
\end{frame}

%%%%%%%%%%%%%%%%%%%%%%%%%%%%%%%%%%%%%%%%%%%%%%%%%%%%%%%%%%%%%%%%%%%%%%%%
\begin{frame}[fragile]\frametitle{Boolean Indexing, Explained}
That last one might be a little strange, so let's make it more clear - \texttt{cities < 1000} returns a Series of \texttt{True/False} values, which we then pass to our Series cities, returning the corresponding \texttt{True} items.
\end{frame}

%%%%%%%%%%%%%%%%%%%%%%%%%%%%%%%%%%%%%%%%%%%%%%%%%%%%%%%%%%%%%%%%%%%%%%%%
\begin{frame}[fragile]\frametitle{Boolean Indexing: Selecting Values}
\begin{lstlisting}
>>> less_than_1000 = cities < 1000
>>> print(less_than_1000)
>>> print('\n')
>>> print(cities[less_than_1000])
Austin            True
Boston           False
Chicago           False
New York          False
Portland          True
San Francisco     False
dtype: bool

Austin      450
Portland    900
dtype: float64
\end{lstlisting}
\end{frame}

%%%%%%%%%%%%%%%%%%%%%%%%%%%%%%%%%%%%%%%%%%%%%%%%%%%%%%%%%%%%%%%%%%%%%%%%
\begin{frame}[fragile]\frametitle{Modifying Series Values by Index}
You can also change the values in a Series on the fly.
\begin{lstlisting}
>>> # changing based on the index
>>> print('Old value:', cities['Chicago'])
>>> cities['Chicago'] = 1400
>>> print('New value:', cities['Chicago'])
Old value: 1000.0
New value: 1400.0
\end{lstlisting}
\end{frame}

%%%%%%%%%%%%%%%%%%%%%%%%%%%%%%%%%%%%%%%%%%%%%%%%%%%%%%%%%%%%%%%%%%%%%%%%
\begin{frame}[fragile]\frametitle{Modifying Values via Boolean Logic}
Changing values using boolean logic
\begin{lstlisting}
>>> print(cities[cities < 1000])
>>> print('\n')
>>> cities[cities < 1000] = 750
>>> print(cities[cities < 1000])
Austin      450
Portland    900
dtype: float64

Austin      750
Portland    750
dtype: float64
\end{lstlisting}
\end{frame}

%%%%%%%%%%%%%%%%%%%%%%%%%%%%%%%%%%%%%%%%%%%%%%%%%%%%%%%%%%%%%%%%%%%%%%%%
\begin{frame}[fragile]\frametitle{Working with Series}
What if you aren't sure whether an item is in the Series? You can check using idiomatic Python.
\begin{lstlisting}
>>> print('Seattle' in cities)
>>> print('San Francisco' in cities)
False
True
\end{lstlisting}
\end{frame}

%%%%%%%%%%%%%%%%%%%%%%%%%%%%%%%%%%%%%%%%%%%%%%%%%%%%%%%%%%%%%%%%%%%%%%%%
\begin{frame}[fragile]\frametitle{Series Arithmetic with Scalars}
Mathematical operations can be done using scalars and functions.
\begin{lstlisting}
>>> # divide city values by 3
>>> cities / 3
Austin           250.000000
Boston                  NaN
Chicago          466.666667
New York         433.333333
Portland         250.000000
San Francisco    366.666667
dtype: float64
\end{lstlisting}
\end{frame}

%%%%%%%%%%%%%%%%%%%%%%%%%%%%%%%%%%%%%%%%%%%%%%%%%%%%%%%%%%%%%%%%%%%%%%%%
\begin{frame}[fragile]\frametitle{Series Arithmetic with NumPy Functions}
\begin{lstlisting}
>>> # square city values
>>> np.square(cities)
Austin            562500
Boston               NaN
Chicago          1960000
New York         1690000
Portland          562500
San Francisco    1210000
dtype: float64
\end{lstlisting}
\end{frame}

%%%%%%%%%%%%%%%%%%%%%%%%%%%%%%%%%%%%%%%%%%%%%%%%%%%%%%%%%%%%%%%%%%%%%%%%
\begin{frame}[fragile]\frametitle{Adding Two Series Together}
You can add two Series together, which returns a union of the two Series with the addition occurring on the shared index values. Values on either Series that did not have a shared index will produce a NULL/NaN (not a number).
\begin{lstlisting}
>>> print(cities[['Chicago', 'New York', 'Portland']])
>>> print('\n')
>>> print(cities[['Austin', 'New York']])
>>> print('\n')
>>> print(cities[['Chicago', 'New York', 'Portland']] + cities[['Austin', 'New York']])
\end{lstlisting}
\end{frame}

%%%%%%%%%%%%%%%%%%%%%%%%%%%%%%%%%%%%%%%%%%%%%%%%%%%%%%%%%%%%%%%%%%%%%%%%
\begin{frame}[fragile]\frametitle{Adding Two Series: Output}
\begin{lstlisting}
Chicago     1400
New York    1300
Portland     750
dtype: float64

Austin       750
New York    1300
dtype: float64

Austin       NaN
Chicago      NaN
New York    2600
Portland     NaN
dtype: float64
\end{lstlisting}
\end{frame}

%%%%%%%%%%%%%%%%%%%%%%%%%%%%%%%%%%%%%%%%%%%%%%%%%%%%%%%%%%%%%%%%%%%%%%%%
\begin{frame}[fragile]\frametitle{Working with Series}
\textbf{NULL Checking}
\begin{itemize}
\item Notice that because Austin, Chicago, and Portland were not found in both Series, they were returned with NULL/NaN values.
\item NULL checking can be performed with \texttt{isnull()} and \texttt{notnull()}.
\end{itemize}
\end{frame}

%%%%%%%%%%%%%%%%%%%%%%%%%%%%%%%%%%%%%%%%%%%%%%%%%%%%%%%%%%%%%%%%%%%%%%%%
\begin{frame}[fragile]\frametitle{Intuition: Index Alignment}
\begin{block}{Intuition}
Adding two Series is like matching up name tags at two different parties and adding values only
where the \emph{same} name appears at both. Anyone who only showed up to one party (an index
value present in just one Series) gets a NaN back: there's nothing on the other side to add
them to.
\end{block}
\end{frame}

%%%%%%%%%%%%%%%%%%%%%%%%%%%%%%%%%%%%%%%%%%%%%%%%%%%%%%%%%%%%%%%%%%%%%%%%
\begin{frame}[fragile]\frametitle{Checking for Non-Null Values}
Return a boolean series indicating which values aren't NULL
\begin{lstlisting}
>>> cities.notnull()

Austin            True
Boston           False
Chicago           True
New York          True
Portland          True
San Francisco     True
dtype: bool
\end{lstlisting}
\end{frame}

%%%%%%%%%%%%%%%%%%%%%%%%%%%%%%%%%%%%%%%%%%%%%%%%%%%%%%%%%%%%%%%%%%%%%%%%
\begin{frame}[fragile]\frametitle{Checking for Null Values}
Using boolean logic to grab the NULL cities
\begin{lstlisting}
>>> print(cities.isnull())
>>> print('\n')
>>> print(cities[cities.isnull()])
Austin           False
Boston            True
Chicago          False
New York         False
Portland         False
San Francisco    False
dtype: bool

Boston   NaN
dtype: float64
\end{lstlisting}
\end{frame}

%%%%%%%%%%%%%%%%%%%%%%%%%%%%%%%%%%%%%%%%%%%%%%%%%%%%%%%%%%
\begin{frame}[fragile]\frametitle{Quick Check: Series}
\begin{block}{Think About It}
\lstinline{cities['Chicago']} returns a single value, but
\lstinline{cities[['Chicago', 'Portland']]} returns a Series. What's the difference in the
syntax, and why does it matter?
\end{block}
\pause
\begin{block}{Answer}
A single string key does scalar/label lookup and returns \emph{one} value. A \emph{list} of
keys is fancy indexing and always returns a Series, even with just one element: the double
brackets are doing two things at once: the outer \lstinline{[]} is the indexing operation, the
inner \lstinline{[]} is a literal Python list being passed as the key.
\end{block}
\end{frame}

%%%%%%%%%%%%%%%%%%%%%%%%%%%%%%%%%%%%%%%%%%%%%%%%%%%%%%%%%%
\begin{frame}
  \begin{center}
    {\Large Dataframes in Pandas}
  \end{center}
\end{frame}

%%%%%%%%%%%%%%%%%%%%%%%%%%%%%%%%%%%%%%%%%%%%%%%%%%%%%%%%%%%
\begin{frame}[fragile]\frametitle{Creating Dataframes}
Creating a DataFrame by passing a numpy array, with a datetime index and labeled columns:
% \begin{center}
% \includegraphics[width=\linewidth,keepaspectratio]{df1}
% \end{center}
\begin{lstlisting}
>>> dates = pd.date_range('20130101', periods=6)
>>> dates
DatetimeIndex(['2013-01-01', '2013-01-02', '2013-01-03', '2013-01-04',
               '2013-01-05', '2013-01-06'],
              dtype='datetime64[ns]', freq='D')

>>> df_demo = pd.DataFrame(np.random.randn(6, 4), index=dates, columns=list('ABCD'))
>>> df_demo
                   A         B         C         D
2013-01-01  0.469112 -0.282863 -1.509059 -1.135632
2013-01-02  1.212112 -0.173215  0.119209 -1.044236
2013-01-03 -0.861849 -2.104569 -0.494929  1.071804
2013-01-04  0.721555 -0.706771 -1.039575  0.271860
2013-01-05 -0.424972  0.567020  0.276232 -1.087401
2013-01-06 -0.673690  0.113648 -1.478427  0.524988
\end{lstlisting}
\end{frame}


%%%%%%%%%%%%%%%%%%%%%%%%%%%%%%%%%%%%%%%%%%%%%%%%%%%%%%%%%%%
\begin{frame}[fragile]\frametitle{Creating Dataframes}
Creating a DataFrame by passing a dict of objects that can be converted to series-like.
% \begin{center}
% \includegraphics[width=\linewidth,keepaspectratio]{df2}
% \end{center}
\begin{lstlisting}
>>> df2 = pd.DataFrame({'A': 1.,
...                      'B': pd.Timestamp('20130102'),
...                      'C': pd.Series(1, index=list(range(4)), dtype='float32'),
...                      'D': np.array([3] * 4, dtype='int32'),
...                      'E': pd.Categorical(["test", "train", "test", "train"]),
...                      'F': 'foo'})
>>> df2
     A          B    C  D      E    F
0  1.0 2013-01-02  1.0  3   test  foo
1  1.0 2013-01-02  1.0  3  train  foo
2  1.0 2013-01-02  1.0  3   test  foo
3  1.0 2013-01-02  1.0  3  train  foo
\end{lstlisting}
\end{frame}

%%%%%%%%%%%%%%%%%%%%%%%%%%%%%%%%%%%%%%%%%%%%%%%%%%%%%%%%%%%
\begin{frame}[fragile]\frametitle{Reading Data from Files}
Every DataFrame so far has been typed in by hand. In practice, almost all data starts life in a file: \texttt{pd.read\_csv()} and \texttt{pd.read\_excel()} load it straight into a DataFrame, inferring column names and types automatically. \texttt{machine\_logs.csv} is provided upfront (no need to type it): copy it into your working folder before class.
\begin{lstlisting}[basicstyle=\tiny\ttfamily]
>>> df = pd.read_csv('machine_logs.csv')
>>> df.head(3)

  machine_id    shift  temperature  vibration operator_notes
0         M1  Morning         72.5      0.021             ok
1         M2  Morning         75.1      0.034  check bearing
2         M1  Evening         78.3      0.041             ok
\end{lstlisting}
\end{frame}

%%%%%%%%%%%%%%%%%%%%%%%%%%%%%%%%%%%%%%%%%%%%%%%%%%%%%%%%%%%
\begin{frame}[fragile]\frametitle{Quick Inspection}
Before doing anything else with a newly loaded DataFrame, check its shape, column types, and non-null counts.
\begin{lstlisting}
>>> df.shape
(6, 5)

>>> df.dtypes
machine_id         object
shift              object
temperature       float64
vibration         float64
operator_notes     object
dtype: object
\end{lstlisting}
\end{frame}

%%%%%%%%%%%%%%%%%%%%%%%%%%%%%%%%%%%%%%%%%%%%%%%%%%%%%%%%%%%
\begin{frame}[fragile]\frametitle{Quick Inspection: Full Structure}
\texttt{df.info()} combines shape, dtypes, and non-null counts into one summary.
\begin{lstlisting}
>>> df.info()
<class 'pandas.core.frame.DataFrame'>
RangeIndex: 6 entries, 0 to 5
Data columns (total 5 columns):
 #   Column          Non-Null Count  Dtype
---  ------          --------------  -----
 0   machine_id      6 non-null      object
 1   shift           6 non-null      object
 2   temperature     6 non-null      float64
 3   vibration       6 non-null      float64
 4   operator_notes  6 non-null      object
dtypes: float64(2), object(3)
\end{lstlisting}
\end{frame}

%%%%%%%%%%%%%%%%%%%%%%%%%%%%%%%%%%%%%%%%%%%%%%%%%%%%%%%%%%%
\begin{frame}[fragile]\frametitle{Viewing Data}
See the top \& bottom rows of the frame
% \begin{center}
% \includegraphics[width=\linewidth,keepaspectratio]{df3}
% \end{center}
\begin{lstlisting}
>>> df_demo.head()
                   A         B         C         D
2013-01-01  0.469112 -0.282863 -1.509059 -1.135632
2013-01-02  1.212112 -0.173215  0.119209 -1.044236
2013-01-03 -0.861849 -2.104569 -0.494929  1.071804
2013-01-04  0.721555 -0.706771 -1.039575  0.271860
2013-01-05 -0.424972  0.567020  0.276232 -1.087401

>>> df_demo.tail(3)
                   A         B         C         D
2013-01-04  0.721555 -0.706771 -1.039575  0.271860
2013-01-05 -0.424972  0.567020  0.276232 -1.087401
2013-01-06 -0.673690  0.113648 -1.478427  0.524988
\end{lstlisting}
\end{frame}

%%%%%%%%%%%%%%%%%%%%%%%%%%%%%%%%%%%%%%%%%%%%%%%%%%%%%%%%%%%
\begin{frame}[fragile]\frametitle{Viewing Data}
Display the index, columns, and the underlying numpy data
% \begin{center}
% \includegraphics[width=\linewidth,keepaspectratio]{df4}
% \end{center}
\begin{lstlisting}
>>> df_demo.index
DatetimeIndex(['2013-01-01', '2013-01-02', '2013-01-03', '2013-01-04',
               '2013-01-05', '2013-01-06'],
              dtype='datetime64[ns]', freq='D')

>>> df_demo.columns
Index(['A', 'B', 'C', 'D'], dtype='object')

>>> df_demo.values
array([[ 0.4691, -0.2829, -1.5091, -1.1356],
       [ 1.2121, -0.1732,  0.1192, -1.0442],
       [-0.8618, -2.1046, -0.4949,  1.0718],
       [ 0.7216, -0.7068, -1.0396,  0.2719],
       [-0.425 ,  0.567 ,  0.2762, -1.0874],
       [-0.6737,  0.1136, -1.4784,  0.525 ]])
\end{lstlisting}
\end{frame}


%%%%%%%%%%%%%%%%%%%%%%%%%%%%%%%%%%%%%%%%%%%%%%%%%%%%%%%%%%%
\begin{frame}[fragile]\frametitle{Viewing Data}
Describe shows a quick statistic summary of your data
% \begin{center}
% \includegraphics[width=\linewidth,keepaspectratio]{df5}
% \end{center}
\begin{lstlisting}
>>> df_demo.describe()
              A         B         C         D
count  6.000000  6.000000  6.000000  6.000000
mean   0.073711 -0.431125 -0.687758 -0.233103
std    0.843157  0.922818  0.779887  0.973118
min   -0.861849 -2.104569 -1.509059 -1.135632
25%   -0.611510 -0.600794 -1.368714 -1.076610
50%    0.022070 -0.228039 -0.767252 -0.386188
75%    0.658444  0.041933 -0.034326  0.461706
max    1.212112  0.567020  0.276232  1.071804
\end{lstlisting}
\end{frame}

%%%%%%%%%%%%%%%%%%%%%%%%%%%%%%%%%%%%%%%%%%%%%%%%%%%%%%%%%%%
\begin{frame}[fragile]\frametitle{Viewing Data}
Sorting by an axis
% \begin{center}
% \includegraphics[width=\linewidth,keepaspectratio]{df6}
% \end{center}
\begin{lstlisting}
>>> df_demo.sort_index(axis=1, ascending=False)
                   D         C         B         A
2013-01-01 -1.135632 -1.509059 -0.282863  0.469112
2013-01-02 -1.044236  0.119209 -0.173215  1.212112
2013-01-03  1.071804 -0.494929 -2.104569 -0.861849
2013-01-04  0.271860 -1.039575 -0.706771  0.721555
2013-01-05 -1.087401  0.276232  0.567020 -0.424972
2013-01-06  0.524988 -1.478427  0.113648 -0.673690
\end{lstlisting}
\end{frame}

%%%%%%%%%%%%%%%%%%%%%%%%%%%%%%%%%%%%%%%%%%%%%%%%%%%%%%%%%%%
\begin{frame}[fragile]\frametitle{Viewing Data}
Sorting by values
% \begin{center}
% \includegraphics[width=\linewidth,keepaspectratio]{df7}
% \end{center}
\begin{lstlisting}
>>> df_demo.sort_values(by='B')
                   A         B         C         D
2013-01-03 -0.861849 -2.104569 -0.494929  1.071804
2013-01-04  0.721555 -0.706771 -1.039575  0.271860
2013-01-01  0.469112 -0.282863 -1.509059 -1.135632
2013-01-02  1.212112 -0.173215  0.119209 -1.044236
2013-01-06 -0.673690  0.113648 -1.478427  0.524988
2013-01-05 -0.424972  0.567020  0.276232 -1.087401
\end{lstlisting}
\end{frame}

%%%%%%%%%%%%%%%%%%%%%%%%%%%%%%%%%%%%%%%%%%%%%%%%%%%%%%%%%%%
\begin{frame}[fragile]\frametitle{Value Counts}
\texttt{value\_counts()} is the fastest way to see how a categorical column is distributed: here, how many readings each machine contributed.
\begin{lstlisting}
>>> df['machine_id'].value_counts()
M1    2
M2    2
M3    2
Name: machine_id, dtype: int64
\end{lstlisting}
\end{frame}

%%%%%%%%%%%%%%%%%%%%%%%%%%%%%%%%%%%%%%%%%%%%%%%%%%%%%%%%%%%
\begin{frame}[fragile]\frametitle{Selection of Data}
Selecting a single column, which yields a Series, equivalent to df\_demo.A
% \begin{center}
% \includegraphics[width=\linewidth,keepaspectratio]{df8}
% \end{center}
\begin{lstlisting}
>>> df_demo['A']
2013-01-01    0.469112
2013-01-02    1.212112
2013-01-03   -0.861849
2013-01-04    0.721555
2013-01-05   -0.424972
2013-01-06   -0.673690
Freq: D, Name: A, dtype: float64
\end{lstlisting}
\end{frame}

%%%%%%%%%%%%%%%%%%%%%%%%%%%%%%%%%%%%%%%%%%%%%%%%%%%%%%%%%%%
\begin{frame}[fragile]\frametitle{Selection of Data}
Selecting via [], which slices the rows.
% \begin{center}
% \includegraphics[width=\linewidth,keepaspectratio]{df9}
% \end{center}
\begin{lstlisting}
>>> df_demo[0:3]
                   A         B         C         D
2013-01-01  0.469112 -0.282863 -1.509059 -1.135632
2013-01-02  1.212112 -0.173215  0.119209 -1.044236
2013-01-03 -0.861849 -2.104569 -0.494929  1.071804

>>> df_demo['20130102':'20130104']
                   A         B         C         D
2013-01-02  1.212112 -0.173215  0.119209 -1.044236
2013-01-03 -0.861849 -2.104569 -0.494929  1.071804
2013-01-04  0.721555 -0.706771 -1.039575  0.271860
\end{lstlisting}
\end{frame}

%%%%%%%%%%%%%%%%%%%%%%%%%%%%%%%%%%%%%%%%%%%%%%%%%%%%%%%%%%%
\begin{frame}[fragile]\frametitle{Selection of Data}
Selecting on a multi-axis by label
% \begin{center}
% \includegraphics[width=\linewidth,keepaspectratio]{df10}
% \end{center}
\begin{lstlisting}
>>> df_demo.loc[:, ['A', 'B']]
                   A         B
2013-01-01  0.469112 -0.282863
2013-01-02  1.212112 -0.173215
2013-01-03 -0.861849 -2.104569
2013-01-04  0.721555 -0.706771
2013-01-05 -0.424972  0.567020
2013-01-06 -0.673690  0.113648
\end{lstlisting}
\end{frame}

%%%%%%%%%%%%%%%%%%%%%%%%%%%%%%%%%%%%%%%%%%%%%%%%%%%%%%%%%%%
\begin{frame}[fragile]\frametitle{Selection of Data}
By integer slices, acting similar to numpy/python
% \begin{center}
% \includegraphics[width=\linewidth,keepaspectratio]{df11}
% \end{center}
\begin{lstlisting}
>>> df_demo.iloc[3:5, 0:2]
                   A         B
2013-01-04  0.721555 -0.706771
2013-01-05 -0.424972  0.567020
\end{lstlisting}
\end{frame}

%%%%%%%%%%%%%%%%%%%%%%%%%%%%%%%%%%%%%%%%%%%%%%%%%%%%%%%%%%%
\begin{frame}[fragile]\frametitle{Selection of Data}
By lists of integer position locations, similar to the numpy/python style
% \begin{center}
% \includegraphics[width=\linewidth,keepaspectratio]{df12}
% \end{center}
\begin{lstlisting}
>>> df_demo.iloc[[1, 2, 4], [0, 2]]
                   A         C
2013-01-02  1.212112  0.119209
2013-01-03 -0.861849 -0.494929
2013-01-05 -0.424972  0.276232
\end{lstlisting}
\end{frame}

%%%%%%%%%%%%%%%%%%%%%%%%%%%%%%%%%%%%%%%%%%%%%%%%%%%%%%%%%%%
\begin{frame}[fragile]\frametitle{Intuition: Label vs. Position}
\begin{block}{Intuition}
Label-based selection is like finding a book by its title on the shelf: you ask for
``Wednesday's row'' or ``column A,'' and pandas looks up that name. Position-based selection is
like grabbing the third book from the left, regardless of what it's called. Both answer valid
but different questions, so pandas gives you separate ways to ask each one instead of guessing
your intent.
\end{block}
\end{frame}

%%%%%%%%%%%%%%%%%%%%%%%%%%%%%%%%%%%%%%%%%%%%%%%%%%%%%%%%%%%
\begin{frame}[fragile]\frametitle{Boolean Indexing}
Using a single column's values to select data.
% \begin{center}
% \includegraphics[width=\linewidth,keepaspectratio]{df13}
% \end{center}
\begin{lstlisting}
>>> df_demo[df_demo.A > 0]
                   A         B         C         D
2013-01-01  0.469112 -0.282863 -1.509059 -1.135632
2013-01-02  1.212112 -0.173215  0.119209 -1.044236
2013-01-04  0.721555 -0.706771 -1.039575  0.271860
\end{lstlisting}
\end{frame}

%%%%%%%%%%%%%%%%%%%%%%%%%%%%%%%%%%%%%%%%%%%%%%%%%%%%%%%%%%%
\begin{frame}[fragile]\frametitle{Setting Values}
Two changes happen here that carry into every slide from this point on: column \texttt{D} is
overwritten entirely, and a new column \texttt{F} is added by assigning a Series that only
covers \emph{some} of the dates: alignment (the same mechanism from the Series section) fills
in the rest with NaN.
\begin{lstlisting}
>>> df_demo.loc[:, 'D'] = np.array([5] * len(df_demo))

>>> s1 = pd.Series([1, 2, 3, 4, 5, 6], index=pd.date_range('20130102', periods=6))
>>> df_demo['F'] = s1
>>> df_demo
              A         B         C  D    F
2013-01-01  0.469112 -0.282863 -1.509059  5  NaN
2013-01-02  1.212112 -0.173215  0.119209  5  1.0
2013-01-03 -0.861849 -2.104569 -0.494929  5  2.0
2013-01-04  0.721555 -0.706771 -1.039575  5  3.0
2013-01-05 -0.424972  0.567020  0.276232  5  4.0
2013-01-06 -0.673690  0.113648 -1.478427  5  5.0
\end{lstlisting}
\end{frame}

%%%%%%%%%%%%%%%%%%%%%%%%%%%%%%%%%%%%%%%%%%%%%%%%%%%%%%%%%%%
\begin{frame}[fragile]\frametitle{Missing Data}
Missing values show up as NaN. Blanking out an earlier value and reindexing to a smaller set
of dates creates some gaps to demonstrate on:
% \begin{center}
% \includegraphics[width=\linewidth,keepaspectratio]{df14}
% \end{center}
\begin{lstlisting}[basicstyle=\tiny\ttfamily]
>>> df_demo.loc[dates[0], ['A', 'B']] = 0
>>> df1 = df_demo.reindex(index=dates[0:4], columns=list(df_demo.columns) + ['E'])
>>> df1.loc[dates[0]:dates[1], 'E'] = 1
>>> df1
                   A         B         C  D    F    E
2013-01-01  0.000000  0.000000 -1.509059  5  NaN  1.0
2013-01-02  1.212112 -0.173215  0.119209  5  1.0  1.0
2013-01-03 -0.861849 -2.104569 -0.494929  5  2.0  NaN
2013-01-04  0.721555 -0.706771 -1.039575  5  3.0  NaN
\end{lstlisting}
\end{frame}

%%%%%%%%%%%%%%%%%%%%%%%%%%%%%%%%%%%%%%%%%%%%%%%%%%%%%%%%%%%
\begin{frame}[fragile]\frametitle{Missing Data: Drop}
To drop any rows that have missing data:
\begin{lstlisting}[basicstyle=\tiny\ttfamily]
>>> df1.dropna(how='any')
                   A         B         C  D    F    E
2013-01-02  1.212112 -0.173215  0.119209  5  1.0  1.0
\end{lstlisting}
\end{frame}

%%%%%%%%%%%%%%%%%%%%%%%%%%%%%%%%%%%%%%%%%%%%%%%%%%%%%%%%%%%
\begin{frame}[fragile]\frametitle{Missing Data: Fill}
Filling missing data:
% \begin{center}
% \includegraphics[width=\linewidth,keepaspectratio]{df15}
% \end{center}
\begin{lstlisting}[basicstyle=\tiny\ttfamily]
>>> df1.fillna(value=5)
                   A         B         C  D    F    E
2013-01-01  0.000000  0.000000 -1.509059  5  5.0  1.0
2013-01-02  1.212112 -0.173215  0.119209  5  1.0  1.0
2013-01-03 -0.861849 -2.104569 -0.494929  5  2.0  5.0
2013-01-04  0.721555 -0.706771 -1.039575  5  3.0  5.0
\end{lstlisting}
\end{frame}

%%%%%%%%%%%%%%%%%%%%%%%%%%%%%%%%%%%%%%%%%%%%%%%%%%%%%%%%%%%
\begin{frame}[fragile]\frametitle{Stats}
Operations in general exclude missing data.
% \begin{center}
% \includegraphics[width=\linewidth,keepaspectratio]{df16}
% \end{center}
Performing a descriptive statistic:
\begin{lstlisting}
>>> df_demo.mean()
A   -0.004474
B   -0.383981
C   -0.687758
D    5.000000
F    3.000000
dtype: float64
\end{lstlisting}
\end{frame}

%%%%%%%%%%%%%%%%%%%%%%%%%%%%%%%%%%%%%%%%%%%%%%%%%%%%%%%%%%%
\begin{frame}[fragile]\frametitle{Stats: The Other Axis}
Same operation on the other axis:
\begin{lstlisting}
>>> df_demo.mean(1)
2013-01-01    0.872735
2013-01-02    1.431621
2013-01-03    0.707731
2013-01-04    1.395042
2013-01-05    1.883656
2013-01-06    1.592306
Freq: D, dtype: float64
\end{lstlisting}
\end{frame}

%%%%%%%%%%%%%%%%%%%%%%%%%%%%%%%%%%%%%%%%%%%%%%%%%%%%%%%%%%%
\begin{frame}[fragile]\frametitle{Apply}
Applying functions to the data
% \begin{center}
% \includegraphics[width=\linewidth,keepaspectratio]{df17}
% \end{center}
\begin{lstlisting}
>>> df_demo.apply(np.cumsum)
                   A         B         C   D     F
2013-01-01  0.000000  0.000000 -1.509059   5   NaN
2013-01-02  1.212112 -0.173215 -1.389850  10   1.0
2013-01-03  0.350263 -2.277784 -1.884779  15   3.0
2013-01-04  1.071818 -2.984555 -2.924354  20   6.0
2013-01-05  0.646846 -2.417535 -2.648122  25  10.0
2013-01-06 -0.026844 -2.303886 -4.126549  30  15.0

>>> df_demo.apply(lambda x: x.max() - x.min())
A    2.073961
B    2.671590
C    1.785291
D    0.000000
F    4.000000
dtype: float64
\end{lstlisting}
\end{frame}

%%%%%%%%%%%%%%%%%%%%%%%%%%%%%%%%%%%%%%%%%%%%%%%%%%%%%%%%%%%
\begin{frame}[fragile]\frametitle{Renaming \& Dropping Columns}
Cleaning up column names and removing ones you don't need is usually the first step of any real analysis.
\begin{lstlisting}
>>> df = df.rename(columns={'vibration': 'vibration_mm_s'})
>>> df = df.drop(columns=['operator_notes'])
>>> df.columns
Index(['machine_id', 'shift', 'temperature', 'vibration_mm_s'], dtype='object')
\end{lstlisting}
\end{frame}

%%%%%%%%%%%%%%%%%%%%%%%%%%%%%%%%%%%%%%%%%%%%%%%%%%%%%%%%%%%
\begin{frame}[fragile]\frametitle{Creating New Columns}
New columns are computed from existing ones with plain vectorized arithmetic and comparisons, no loops needed.
\begin{lstlisting}
>>> df['vibration_in'] = df['vibration_mm_s'] / 25.4
>>> df['high_vibration'] = df['vibration_mm_s'] > 0.025
>>> df[['machine_id', 'vibration_mm_s', 'high_vibration']]

  machine_id  vibration_mm_s  high_vibration
0         M1           0.021           False
1         M2           0.034            True
2         M1           0.041            True
3         M2           0.028            True
4         M3           0.019           False
5         M3           0.052            True
\end{lstlisting}
\end{frame}

%%%%%%%%%%%%%%%%%%%%%%%%%%%%%%%%%%%%%%%%%%%%%%%%%%%%%%%%%%%
\begin{frame}[fragile]\frametitle{Quick Check: DataFrames}
\begin{block}{Think About It}
\lstinline{df['A']} returns a Series (a column), but \lstinline{df[0:2]} returns a DataFrame
(the first two rows). Same \lstinline{[]} operator, same object: why does pandas treat a
string key and a slice so differently?
\end{block}
\pause
\begin{block}{Answer}
The plain \lstinline{[]} operator on a DataFrame is intentionally overloaded: a single label
looks up a \emph{column} (matching the \lstinline{df.A} shortcut), while a slice looks up
\emph{rows}, mirroring plain Python/NumPy slicing. Pandas infers your intent from the
\emph{type} of the key rather than a fixed row/column rule, which is exactly why
\lstinline{.loc}/\lstinline{.iloc} exist: they remove the ambiguity by always meaning
``by label'' or ``by position,'' regardless of what kind of key you pass.
\end{block}
\end{frame}

%%%%%%%%%%%%%%%%%%%%%%%%%%%%%%%%%%%%%%%%%%%%%%%%%%%%%%%%%%%
\begin{frame}[fragile]\frametitle{GroupBy: Split-Apply-Combine}
\texttt{groupby()} splits rows into groups by a key column. Nothing is computed until you chain on an aggregate.
\begin{block}{Intuition}
Think of \texttt{groupby()} as sorting mail into pigeonholes by recipient. Nothing gets read or
totalled yet: that's just the sorting step. The real work happens when a clerk (an aggregate
like \texttt{.mean()} or \texttt{.sum()}) processes each pigeonhole separately and the results
are stacked back into a single tray.
\end{block}
\begin{lstlisting}
>>> df.groupby('machine_id')
<pandas.core.groupby.generic.DataFrameGroupBy object at 0x...>
\end{lstlisting}
\end{frame}

%%%%%%%%%%%%%%%%%%%%%%%%%%%%%%%%%%%%%%%%%%%%%%%%%%%%%%%%%%%
\begin{frame}[fragile]\frametitle{GroupBy with Aggregation}
Chain \texttt{.mean()}, \texttt{.sum()}, or \texttt{.agg()} to actually run the per-group computation.
\begin{lstlisting}
>>> df.groupby('machine_id')['temperature'].mean()
M1    75.40
M2    74.55
M3    75.50
Name: temperature, dtype: float64

>>> df.groupby('machine_id').agg({'temperature': 'mean', 'vibration_mm_s': 'max'})
            temperature  vibration_mm_s
machine_id
M1                75.40            0.041
M2                74.55            0.034
M3                75.50            0.052
\end{lstlisting}
\end{frame}

%%%%%%%%%%%%%%%%%%%%%%%%%%%%%%%%%%%%%%%%%%%%%%%%%%%%%%%%%%%
\begin{frame}[fragile]\frametitle{Quick Check: GroupBy}
\begin{block}{Think About It}
\lstinline{df.groupby('machine_id')} on its own prints only an object description, not a table
of numbers. Why doesn't pandas compute anything at that point?
\end{block}
\pause
\begin{block}{Answer}
\lstinline{groupby()} only performs the \emph{split}: it doesn't yet know what you want to
\emph{apply} to each group. Waiting until you chain on \lstinline{.mean()}, \lstinline{.sum()},
or \lstinline{.agg()} means the same grouped object can be reused for several different
aggregations without repeating the split step, and pandas never computes a result you never
asked for.
\end{block}
\end{frame}

%%%%%%%%%%%%%%%%%%%%%%%%%%%%%%%%%%%%%%%%%%%%%%%%%%%%%%%%%%%
\begin{frame}[fragile]\frametitle{Concatenating DataFrames}
\texttt{pd.concat()} stacks DataFrames with the same columns on top of each other, useful when the same kind of data arrives in separate batches (e.g., one file per shift).
\begin{lstlisting}
>>> morning = df[df['shift'] == 'Morning'][['machine_id', 'temperature']]
>>> evening = df[df['shift'] == 'Evening'][['machine_id', 'temperature']]
>>> pd.concat([morning, evening])
  machine_id  temperature
0         M1         72.5
1         M2         75.1
4         M3         69.8
2         M1         78.3
3         M2         74.0
5         M3         81.2
\end{lstlisting}
\end{frame}

%%%%%%%%%%%%%%%%%%%%%%%%%%%%%%%%%%%%%%%%%%%%%%%%%%%%%%%%%%%
\begin{frame}[fragile]\frametitle{Merging DataFrames}
\texttt{pd.merge()} joins two DataFrames on a shared key column: like a VLOOKUP, or a SQL JOIN.
\begin{lstlisting}
>>> machine_info = pd.DataFrame({
...     'machine_id': ['M1', 'M2', 'M3'],
...     'machine_type': ['Lathe', 'Mill', 'Drill']
... })
>>> pd.merge(df, machine_info, on='machine_id')[
...     ['machine_id', 'machine_type', 'temperature']].head(3)
  machine_id machine_type  temperature
0         M1        Lathe         72.5
1         M2         Mill         75.1
2         M1        Lathe         78.3
\end{lstlisting}
\end{frame}

%%%%%%%%%%%%%%%%%%%%%%%%%%%%%%%%%%%%%%%%%%%%%%%%%%%%%%%%%%%
\begin{frame}[fragile]\frametitle{Plotting with Pandas}
\texttt{.plot()} is a thin wrapper around the \texttt{matplotlib} imported all the way back on the very first Pandas slide, no separate charting library needed for a first look at the data.
\begin{lstlisting}
>>> avg_temp = df.groupby('machine_id')['temperature'].mean()
>>> avg_temp.plot(kind='bar')
>>> plt.ylabel('Avg Temperature (deg C)')
>>> plt.show()
\end{lstlisting}
\end{frame}

%%%%%%%%%%%%%%%%%%%%%%%%%%%%%%%%%%%%%%%%%%%%%%%%%%%%%%%%%%%
\begin{frame}[fragile]\frametitle{Putting It Together: Machine Health Check}
The same five-step pattern (read, inspect, clean, transform, summarize) shows up in almost every dataset you'll touch in this course.
\begin{lstlisting}
>>> df = pd.read_csv('machine_logs.csv')
>>> df = df.rename(columns={'vibration': 'vibration_mm_s'})
>>> df = df.drop(columns=['operator_notes'])
>>> df['high_vibration'] = df['vibration_mm_s'] > 0.025

>>> summary = df.groupby('machine_id').agg(
...     avg_temp=('temperature', 'mean'),
...     flagged_readings=('high_vibration', 'sum')
... )
>>> summary
            avg_temp  flagged_readings
machine_id
M1             75.40                 1
M2             74.55                 2
M3             75.50                 1
\end{lstlisting}
\end{frame}

%%%%%%%%%%%%%%%%%%%%%%%%%%%%%%%%%%%%%%%%%%%%%%%%%%%%%%%%%%%
\begin{frame}[fragile]\frametitle{Putting It Together: Visualizing the Summary}
One line turns the summary table into a chart flagging which machine needs a closer look.
\begin{lstlisting}
>>> summary['avg_temp'].plot(kind='bar', color='steelblue')
>>> plt.title('Average Temperature by Machine')
>>> plt.ylabel('deg C')
>>> plt.show()
\end{lstlisting}
\end{frame}

%%%%%%%%%%%%%%%%%%%%%%%%%%%%%%%%%%%%%%%%%%%%%%%%%%%%%%%%%%%
\begin{frame}[fragile]\frametitle{Try It Yourself}
Using the same \texttt{machine\_logs.csv} you already have, no new data or setup needed:
\begin{itemize}
\item Which machine has the highest average vibration?
\item Filter the DataFrame to only the Evening shift rows.
\item Add a column flagging rows where \texttt{temperature > 76}.
\end{itemize}
\end{frame}
